\begin{thebibliography}{99}
\small\setlength{\itemsep}{6pt}
\bibitem{Repository} I. Dorofeev.
\emph{Polynomial Rectification of the Shastri Trefoil and a Harmonic
Figure-Eight}. Рабочий препринт и вычислительные сертификаты.
\href{https://github.com/qft98romanov-hash/abhyankar-moh-trefoil/tree/b954f3609e900761fddc9c8ff03d8301e9604cab}{Репозиторий, проверенная версия}.

\bibitem{AM} S. S. Abhyankar, T. T. Moh.
\emph{Embeddings of the line in the plane}.
J. reine angew. Math. \textbf{276} (1975), 148--166.
\href{https://doi.org/10.1515/crll.1975.276.148}{DOI: 10.1515/crll.1975.276.148}.

\bibitem{Suzuki} M. Suzuki.
\emph{Propriétés topologiques des polynômes de deux variables complexes,
et automorphismes algébriques de l'espace $\C^2$}.
J. Math. Soc. Japan \textbf{26} (1974), 241--257.
\href{https://doi.org/10.2969/jmsj/02620241}{DOI: 10.2969/jmsj/02620241}.

\bibitem{Shastri} A. R. Shastri.
\emph{Polynomial representations of knots}.
Tôhoku Math. J. (2) \textbf{44} (1992), 11--17.
\href{https://doi.org/10.2748/tmj/1178227371}{DOI: 10.2748/tmj/1178227371};
\href{https://www.jstage.jst.go.jp/article/tmj1949/44/1/44_1_11/_pdf}{открытый PDF}.
Вопросы Q.1--Q.2 и Theorem 1, с.~12; Example 1, с.~15.

\bibitem{Craighero} P. C. Craighero.
\emph{About Abhyankar's conjectures on space lines}.
Rend. Sem. Mat. Univ. Padova \textbf{74} (1985), 115--122.
\href{https://www.numdam.org/item/RSMUP_1985__74__115_0.pdf}{Открытый PDF}.
Исторические формулировки: введение, с.~115--116.

\bibitem{vdER} A. van den Essen, P. van Rossum.
\emph{Triangular derivations related to problems on affine $n$-space}.
Proc. Amer. Math. Soc. \textbf{130} (2002), no.~5, 1311--1322.
\href{https://doi.org/10.1090/S0002-9939-01-06212-8}{DOI: 10.1090/S0002-9939-01-06212-8}.
Conjectures 5.1--5.3, Proposition 5.4, Corollary 5.5, с.~1319--1321.

\bibitem{Masuda} K. Masuda.
\emph{Equivariant embeddings and semi-invariant locally nilpotent
derivations}. Авторская версия.
\href{https://sci-tech.ksc.kwansei.ac.jp/~kayo/cmn/papers/eembed.pdf}{Открытый PDF}.
Theorems 1.1--1.2; градуировка является существенной частью постановки.

\bibitem{EVM} E. Edo, A. van den Essen, S. Maubach.
\emph{A note on $k[z]$-automorphisms in two variables}.
J. Pure Appl. Algebra \textbf{213} (2009), 1197--1200.
\href{https://doi.org/10.1016/j.jpaa.2008.11.040}{DOI: 10.1016/j.jpaa.2008.11.040};
\href{https://arxiv.org/abs/0809.0767}{arXiv:0809.0767}.
Proposition 3.1 в открытой версии.

\bibitem{FF} G. Freudenburg, J. Freudenburg.
\emph{Curves defined by Chebyshev polynomials} (2009).
\href{https://arxiv.org/abs/0902.3440}{arXiv:0902.3440}.

\bibitem{KP} P.-V. Koseleff, D. Pecker.
\emph{Harmonic knots}.
J. Knot Theory Ramifications \textbf{25} (2016), 1650074.
\href{https://doi.org/10.1142/S0218216516500747}{DOI: 10.1142/S0218216516500747};
\href{https://arxiv.org/abs/1203.4376}{arXiv:1203.4376}.
Семейство $T(2,2n+1)$ и рисунок 3 для $H(4,5,7)$.

\bibitem{Chau} N. V. Chau.
\emph{A simple proof of Jung's theorem on polynomial automorphisms
of $\C^2$} (2004).
\href{https://arxiv.org/abs/math/0408077}{arXiv:math/0408077}.

\bibitem{DG} N. Dasgupta, S. A. Gaifullin.
\emph{On locally nilpotent derivations of polynomial algebra in three
variables}. Sbornik: Mathematics \textbf{216}:4 (2025), 456--484.
\href{https://doi.org/10.4213/sm10094e}{DOI: 10.4213/sm10094e}.
Русская версия: Математический сборник \textbf{216}:4 (2025), 3--34,
\href{https://doi.org/10.4213/sm10094}{DOI: 10.4213/sm10094}.
Theorem 4.12, Example 4.19, Remark 4.20, Question 4.21.

\bibitem{FKZ} H. Flenner, S. Kaliman, M. Zaidenberg.
\emph{A Gromov--Winkelmann type theorem for flexible varieties}.
J. Eur. Math. Soc. \textbf{18} (2016), 2483--2510.
\href{https://doi.org/10.4171/JEMS/646}{DOI: 10.4171/JEMS/646};
\href{https://ems.press/content/serial-article-files/32154}{открытый PDF}.

\bibitem{Nowicki} A. Nowicki.
\emph{The divergence of polynomial derivations}.
Авторский текст от 29 декабря 1993 года, Theorem 3.3.
\href{https://www.researchgate.net/publication/242147525_The_divergence_of_polynomial_derivations}{Открытая авторская копия}.

\bibitem{Freudenburg} G. Freudenburg.
\emph{Algebraic Theory of Locally Nilpotent Derivations}.
2nd enlarged edition, Encyclopaedia of Mathematical Sciences,
Springer, 2017.
\href{https://doi.org/10.1007/978-3-662-55350-3}{DOI: 10.1007/978-3-662-55350-3}.
\end{thebibliography}
